\documentclass[11pt,a4paper]{article}
\usepackage{fontspec,amsmath,amssymb,amsthm,microtype}
\setmainfont{lmroman10-regular.otf}[BoldFont=lmroman10-bold.otf,ItalicFont=lmroman10-italic.otf,BoldItalicFont=lmroman10-bolditalic.otf]
\setmonofont{lmmono10-regular.otf}[Scale=MatchLowercase]
\usepackage[margin=25mm]{geometry}
\usepackage[unicode,hidelinks]{hyperref}
\setlength{\emergencystretch}{3em}
\setlength{\parskip}{0.4em}
\setlength{\parindent}{0pt}
\allowdisplaybreaks[2]
\raggedbottom
\title{\texorpdfstring{A uniform cubic remainder for smooth test functions}{A uniform cubic remainder for smooth test functions}}
\author{LangevinLean mathematical exposition}
\date{}
\begin{document}\maketitle
\noindent Source module: \texttt{Unadjusted\allowbreak{}Langevin/\allowbreak{}Weak\allowbreak{}Forward\allowbreak{}Taylor.\allowbreak{}lean}. The accompanying HTML pairs this mathematical exposition with the exact Lean statements and records the scope of the formalization.


Second-order expansions of smooth compactly supported tests have a remainder bounded by a constant times \(\|v\|^3\). These analytic estimates support weak forward-equation calculations.


\section{\texorpdfstring{second\allowbreak{}Order\allowbreak{}Expansion}{secondOrderExpansion}}


Define \(T_2\varphi(x;v)=\varphi(x)+D\varphi(x)[v]+\frac12D^2\varphi(x)[v,v]\). The iterated Fréchet derivative is a continuous bilinear map, evaluated twice at the same displacement. This separates the deterministic polynomial from the Gaussian expectation calculations performed later.



{\raggedright\noindent Lean declarations: \texttt{second\allowbreak{}Order\allowbreak{}Expansion}.\par}

\section{\texorpdfstring{second\allowbreak{}Order\allowbreak{}Expansion remainder le}{secondOrderExpansion remainder le}}


Restrict \(\varphi\) to the segment \(t\mapsto x+tv\). Its third scalar derivative is \(D^3\varphi(x+tv)[v,v,v]\), bounded by \(C\|v\|^3\). Taylor's theorem with the factorial remainder gives \(|\varphi(x+v)-T_2\varphi(x;v)|\le(C/6)\|v\|^3\). Lean identifies the first three coefficients with the iterated Fréchet derivatives.



{\raggedright\noindent Lean declarations: \texttt{second\allowbreak{}Order\allowbreak{}Expansion\_\allowbreak{}remainder\_\allowbreak{}le}.\par}

\section{\texorpdfstring{compact\allowbreak{}Smooth second\allowbreak{}Order\allowbreak{}Expansion remainder}{compactSmooth secondOrderExpansion remainder}}


Smoothness makes \(D^3\varphi\) continuous, and compact support of \(\varphi\) implies compact support of its derivatives. Hence its operator norm is globally bounded by some nonnegative constant. Apply the previous Taylor theorem and absorb the factor \(1/6\) into that constant to obtain a uniform cubic remainder bound.



{\raggedright\noindent Lean declarations: \texttt{compact\allowbreak{}Smooth\_\allowbreak{}second\allowbreak{}Order\allowbreak{}Expansion\_\allowbreak{}remainder}.\par}

\end{document}
